\originalchapter{研究生作业九：场变换与相对论轨道}
题面据8.311作业9全部6题. 解答为AI独立推导并参照官方解答. 采用上一章被动boost约定, 原题源变换的加号对应相反boost方向, 不与本章场式混用. 原卷用$\beta$表示Lorentz因子, 本书改用$\gamma$, $\beta=v/c$.
\originalsection{带电直线的变换}
\begin{exercise}
原题1. (a) 静止带电无限直线的密度经沿线运动系如何改变? 求电流并核四流; (b) 两系求场并核场变换.
\end{exercise}
\begin{solution}
静止线荷$\lambda$, $I=0$. 长度收缩而电荷不变给$\lambda'=\gamma\lambda$, 电荷在$K'$沿$-z$运动给$I'=-v\lambda'$. 四流同boost给$\rho'=\gamma(\rho-vJ_z/c^2)$、$J'_z=\gamma(J_z-v\rho)$, 正好一致. 静止$E_s=\lambda/(2\pi\eps s)$、$B_\phi=0$. 运动系高斯/安培给$E'_s=\gamma E_s$, $B'_\phi=\mu_0 I'/(2\pi s)=-\gamma vE_s/c^2$. 沿$+z$被动boost的横场式$\ee'=\gamma(\ee+\vct v\times\bb)$、$\bb'=\gamma(\bb-\vct v\times\ee/c^2)$也给同值. 原卷源加号与场减号若同用会使电流磁场方向冲突.
\end{solution}
\originalsection{运动介质板与导体极限}
\begin{exercise}
原题2. 板以$\vct v$运动, 实验室外场$\ee=0$, $\bb=\bb_0\perp\vct v$, 两者平行板面. (a) 共动系求内外场、极化和面荷, 变回实验室; (b) 理想导体重做.
\end{exercise}
\begin{solution}
板静止系（撇号）外场$\ee'_{\rm out}=\gamma\vct v\times\bb_0$, $\bb'_{\rm out}=\gamma\bb_0$. 非磁介质$H'_t$连续、无自由面荷$D'_n$连续给$\ee'_{\rm in}=\ee'_{\rm out}/\varepsilon_r$、$\bb'_{\rm in}=\gamma\bb_0$. $\vct P'=\eps(\varepsilon_r-1)\ee'_{\rm in}$, 上下表面$\sigma_b'=\pm|P'|$, 法线按$\vct v\times\bb_0$正方向.

反变换给$\ee_{\rm in}=-\gamma^2(1-1/\varepsilon_r)\vct v\times\bb_0$, $\bb_{\rm in}=\gamma^2(1-\beta^2/\varepsilon_r)\bb_0$, 外场仍$0,\bb_0$. 面荷沿面运动, $\sigma_b=\gamma\sigma_b'=\pm\eps\gamma^2(1-1/\varepsilon_r)|\vct v\times\bb_0|$, 并有面流$\vct K=\sigma_b\vct v$. 其电场跳变与磁场跳变恰好还原上述内场. 这是无限平板、外场由固定源维持的线性模型.

导体在自身系$\ee'_{\rm in}=0$, 对无独立磁屏蔽背景取$\bb'_{\rm in}=\gamma\bb_0$. 因而实验室$\ee_{\rm in}=-\gamma^2\vct v\times\bb_0$, $\bb_{\rm in}=\gamma^2\bb_0$, $\sigma=\pm\eps\gamma^2vB_0$. 形式上是$\varepsilon_r\to\infty$极限, 小速下恢复$\ee_{\rm in}=-\vct v\times\bb_0$; 理想导体不等于超导磁通排除.
\end{solution}
\originalsection{Compton能量动量消元}
\begin{exercise}
原题3. 光子波长$\lambda$撞静止电子, 散射角$\theta$、波长$\lambda'$. (a) 写相对论守恒; (b) 求波长差.
\end{exercise}
\begin{solution}
令光子能量$\epsilon=hc/\lambda$, 散射$\epsilon'=hc/\lambda'$, 电子末能$E_e=mc^2+\epsilon-\epsilon'$. 动量$\pp_e=(\epsilon\nn_i-\epsilon'\nn_f)/c$, 故$p_e^2c^2=\epsilon^2+\epsilon'^2-2\epsilon\epsilon'\cos\theta$. 代入$E_e^2=m^2c^4+p_e^2c^2$, 消去平方项得$mc^2(\epsilon-\epsilon')=\epsilon\epsilon'(1-\cos\theta)$, 即$\lambda'-\lambda=h(1-\cos\theta)/(mc)\ge0$. 零角零位移、反向为两倍Compton波长是两项核对.
\end{solution}
\originalsection{恒磁场、交叉场及类光例外}
\begin{exercise}
原题4. $\dot{\pp}=q(\ee+\vct v\times\bb)$, $\vct v=c^2\pp/\mathcal E$. (a) 恒磁场求螺旋; (b) 匀强相互正交$\ee,\bb$求可化纯场系及轨迹, 并检查原卷“总能消去一场”的边界.
\end{exercise}
\begin{solution}
(a) 磁力不做功, $\gamma$常数, 沿$B\hat{\vct z}$的$p_z,v_z$恒定. 横方程$\dot p_x=\Omega p_y$, $\dot p_y=-\Omega p_x$, $\Omega=qB/(\gamma m)$. 因而$p_x=p_{x0}\cos\Omega t+p_{y0}\sin\Omega t$, $p_y=p_{y0}\cos\Omega t-p_{x0}\sin\Omega t$. 积分给$x=x_0+[p_{x0}\sin\Omega t+p_{y0}(1-\cos\Omega t)]/(\gamma m\Omega)$、$y=y_0+[p_{y0}\sin\Omega t-p_{x0}(1-\cos\Omega t)]/(\gamma m\Omega)$及$z=z_0+v_zt$. 半径$p_\perp/(|q|B)$, 频率为非相对论频率的$1/\gamma$.

(b) 若$E<cB$, 选$\vct w=\ee\times\bb/B^2$使$\ee'=0$, $|w|=E/B<c$. 在新系用(a)螺旋再反boost得到实验室参数轨迹. 若$E>cB$, 选$\vct w=c^2\ee\times\bb/E^2$使$\bb'=0$. 取新系$E'$沿$x'$, $p'_x=p'_{x0}+qE't'$, $p'_{y,z}$不变; 令$M_\perp^2=m^2+p_\perp'^2/c^2$, 则$x'-x'_0=[\mathcal E'(t')-\mathcal E'_0]/(qE')$, $y'-y'_0=(cp'_y/qE')[\operatorname{arsinh}(p'_x/(M_\perp c))]_{0}^{t'}$, $z'$同式. 这些连同$t'$反boost是完整轨迹参数式. 具体令$\hat{\vct w}$为boost方向, 初四动量变换为$\mathcal E'_0=\gamma_w(\mathcal E_0-\vct w\cdot\pp_0)$、$p'_{0\parallel}=\gamma_w(p_{0\parallel}-w\mathcal E_0/c^2)$、$\pp'_{0\perp}=\pp_{0\perp}$. 初位置按同一被动boost变换. 将这些初值代入纯场解后, 逆变换$t=\gamma_w(t'+\vct w\cdot\rr'/c^2)$、$r_\parallel=\gamma_w(r'_\parallel+wt')$、$\rr_\perp=\rr'_\perp$, 即闭合地给实验室轨迹.

非零$E=cB$的两个候选$w$均达$c$, 不存在消去一场的惯性系. 为补此原题遗漏, 取$\ee=E\hat{\vct x}$、$\bb=-E\hat{\vct z}/c$, 无量纲四速度$u^0=\gamma$, $u^i=\gamma v_i/c$, $a=qE/(mc)$. 固有时方程$\dot u^0=a u^x$, $\dot u^x=a(u^0-u^y)$、$\dot u^y=a u^x$使$D=u^0-u^y$恒定. 解为$u^x=u^x_0+aD\tau$, $u^0=u^0_0+au^x_0\tau+a^2D\tau^2/2$, $u^y=u^y_0+au^x_0\tau+a^2D\tau^2/2$, $u^z=u^z_0$. 再积分$\dd t/\dd\tau=u^0$、$\dd\rr/\dd\tau=c\vct u$即得到$x$二次、$t,y$三次的显式多项式轨迹. 这是独立补充的类光支, 不冒称原官方解已包含它.
\end{solution}
\originalsection{相对论Kepler与近日点推进}
\begin{exercise}
原题5. $L=-mc^2\sqrt{1-(\dot\rho^2+\rho^2\dot\phi^2)/c^2}+\kappa/\rho$. 求$M,E$及径向一阶式; (a) 圆轨道周期; (b) 一般非圆轨道及有界支进动、低速推进量.
\end{exercise}
\begin{solution}
无显$\phi,t$给$M=\gamma m\rho^2\dot\phi$, $E=\gamma mc^2-\kappa/\rho$. $p_\rho=\gamma m\dot\rho$满足$p_\rho^2=(E+\kappa/\rho)^2/c^2-m^2c^2-M^2/\rho^2$, 所以$\dot\rho=c^2p_\rho/(E+\kappa/\rho)$, 正负支由初值选.

(a) 圆轨道$\gamma mv^2\rho=\kappa$, $T^2=4\pi^2\gamma m\rho^3/\kappa$. 这里$\gamma$随$\rho$隐变, 令$b=\kappa/(mc^2\rho)$, 有$\gamma-1/\gamma=b$, 即$\gamma=(b+\sqrt{b^2+4})/2$. 不可把$\gamma$当所有圆轨道相同的常数.

(b) $u=1/\rho$有$u'=-p_\rho/M$, 因而$u''+\eta^2u=E\kappa/(c^2M^2)$, $\eta^2=1-\kappa^2/(c^2M^2)$. $|M|>\kappa/c$的有界支为$u=u_0[1+e\cos\eta(\phi-\phi_0)]$, $u_0=E\kappa/(c^2M^2-\kappa^2)$. 有界非碰撞还需$0<E<mc^2$和$e<1$, 其中$e^2=1+(E^2-m^2c^4)(c^2M^2-\kappa^2)/(E^2\kappa^2)$. 一径向周期角增量$2\pi/\eta$, 近日点推进$\Delta\phi=2\pi(1/\eta-1)\simeq\pi\kappa^2/(c^2M^2)$. 低速可写$\pi v_{\rm circ}^2/c^2$的近圆尺度. 一般非圆解需分支: $\eta^2>0$为$u=C/\eta^2+A\cos\eta\phi+B\sin\eta\phi$, $C=E\kappa/(c^2M^2)$; 非有界散射同样取此解中物理$u>0$区间. $\eta=0$为$u=C\phi^2/2+A\phi+B$; $\eta^2=-\chi^2<0$为$u=-C/\chi^2+Ae^{\chi\phi}+Be^{-\chi\phi}$. 各支均须$u>0$、$E+\kappa u>0$, 常量由径向初值确定. $M=0$单独用径向一阶式, 可能塌心, 不称进动椭圆. 有界正旋向支周期频率满足$\Omega_\phi/\Omega_r=1/\eta$, 平均近日点率$(1/\eta-1)\Omega_r$, 反旋向改有向符号.
\end{solution}
\originalsection{火箭质量比与快度}
\begin{exercise}
原题6. 初静止火箭质量$M_0$, 排气在瞬时共动系速度$-u$, $0<u<c$. (a) Newton质量比; (b) 相对论推广并核低速极限.
\end{exercise}
\begin{solution}
(a) 短时动量守恒$M\dd v=-u\dd M$, 积分$v=u\ln(M_0/M)$, 即$M/M_0=e^{-v/u}$. (b) 用剩余火箭静质量$M$与快度$\eta$. 瞬时静止系排出物的总能量$dE_{\rm out}=-c^2\dd M$至一阶, 排出动量$-u\,dE_{\rm out}/c^2$. 火箭反冲$Mc\dd\eta=-u\dd M$, 得$\eta=(u/c)\ln(M_0/M)$, $v=c\tanh\eta$, 故
$M/M_0=\exp[-(c/u)\operatorname{artanh}(v/c)]=[(1-v/c)/(1+v/c)]^{c/(2u)}$.
排出物静质量不直接等于$-\dd M$, 还含其相对论能量及燃料内能. 当$v/c\ll1$时恢复Newton式; 单凭$u\ll c$却允许极大质量比, 仍可能达到相对论$v$, 所以原低速极限须同时限定最终$v$.
\end{solution}
